\section{Introduction}
\label{sec:geometric-results}
Near the first time at which a prescribed number of collisions is possible,
a dispersing billiard follows a nearly normal alternating channel. The
probability of such a record decreases exponentially with its length.
Its stationary action, however, has only two free endpoints. This raises
a relative approximation problem: can the nonlinear physical law be
controlled on a contact neighborhood which does not decrease as the
number of collisions increases? A positive answer would produce an
observable boundary invariant beyond the gap and the linear multiplier:
one determined by physical collision probabilities rather than by an
assumed access to endpoint actions.

The distinction between absolute and relative control is essential.
The endpoint Hessian stays uniformly positive, while the mixed endpoint
derivative, which supplies the physical flux, tends exponentially to
zero. Dividing an absolute error estimate by this derivative would lose
precisely the uniformity needed to compare rare events. We instead
normalize the cofactor formula before passing to the limit. The interior
Hessian perturbation is localized near the two ends in trace norm. Its
relative determinant consequently separates into two half-line factors,
with no factor proportional to the number of interior collisions.

Our first result is a nonlinear boundary law on a common closed offset
collar, with every fixed geometric and offset derivative. It applies
locally about any separated, positively curved, smooth periodic
configuration; neither symmetry nor finite horizon is required. The
limit retains the physical residual-time constraint. It is not a product
probability, even though its amplitude is a product of two boundary
factors. The full statement is given below, after the observation has
been specified.

The two oriented even subsequences determine an invariant of the two
contacts. The odd subsequence carries an independent compatibility
constraint. From the two same-type scalar laws we recover the complete
normalized boundary energy profiles, without assuming identical or
even facing graphs. In the norm of three offset derivatives this
inverse is locally Lipschitz. The mixed-type law is then determined
by a Volterra convolution. Its transformed coefficient matrix has
rank one to every finite order; for example, its second derivative
is the mean of the two same-type second derivatives minus an explicit
nonnegative square. The finite-record defect in these identities is
exponentially small in the bridge length. These are nonlinear
consequences of the long-bridge theorem, not consequences of a
finite-dimensional fit to leading amplitudes.

The profile has a direct geometric meaning: it is the normalized
pushforward of a half-line determinant-weighted endpoint measure by
its excess action. The two branches of the contact coordinate are
summed, not individually recovered. The nonlinear realizations below
vary this measure while fixing the selected leading data. Thus the
invariant distinguishes information invisible to the gap and multiplier,
without asserting rigidity of an arbitrary asymmetric table.

We also recover these full profiles from finite physical preparations.
Theorem~\ref{thm:v10-acquisition} requires only selected-event binary
reports on a positive-offset grid, at a common even flight number.
A regularized minimum-discrepancy procedure controls scalar noise,
finite-bridge bias and derivative approximation separately. For profiles
in a bounded $C^{m-1,1}$ class, $m\ge4$, Corollary~\ref{cor:v10-budget}
gives the sufficient preparation count
\[
 C\varepsilon^{-\left(2+8/(m-3)+\gamma/|\log\tau|\right)}
                         \log\frac{C}{\eta\varepsilon}
\]
for uniform accuracy $\varepsilon$ at confidence $1-\eta$.
Here $\tau$ is a supplied relative convergence rate and the onset,
area, multiplier and patch labels are known. All failed preparations
are charged. This theorem turns the differentiated-data inverse into
a finite-observation statement for the same boundary invariant. Its
cost is an upper bound, not a claimed optimal rate or a consequence
of the separate unlabelled parametric experiment.

Our second group of results compares experiments derived from this law.
For endpoint and residual-time observations, relative control passes to
total variation while preserving the rare-event mass. A sharp quadratic
support calculation and a uniform tangent approximation then identify a
critical testing profile, including the supercritical regime. This last
profile is a joint small-offset limit. It does not identify a sharp
fixed-positive-offset rate for the nonlinear boundary approximation.

For unlabelled binary counts, a different issue arises: contact
curvatures may coalesce. In a specified analytic physical family, four
fixed physical windows recover the unknown gap and the symmetric
curvature coordinates. If curvature accuracy is $\varepsilon$ and gap
accuracy is $\delta$, the optimal worst-case expected preparation order,
among bounded-flight binary designs in a fixed local collar, is
\[
                (\varepsilon^{-6}+\delta^{-2})\log(1/\eta)
\]
at confidence $1-\eta$. The upper bound has a deterministic preparation
count. Its two lower bounds come from distinct physical alternatives:
symmetric curvature splitting and a displacement of the onset. Thus a
curvature-only comparison and a comparison with finer timing accuracy
are genuinely different optimization problems.

Fixed-window inversion uses the exact analytic family. We also give an
explicit smooth remainder experiment in which the higher offset
coefficients are unknown. A count-only procedure retains the leading
inverse, acquires its own bracket, and estimates the gap to order
$h^{m+1}$ and the curvature multiset to order $h^{m/3}$ with cost
$h^{-(2m+2)}\log(1/\eta)$ for each fixed $m$. No exact higher-order
probability formula is supplied in that experiment. With a strict
interior margin on the nuisance bounds, its optimal joint cost is
$(\varepsilon^{-(6+6/m)}+\delta^{-2})\log(1/\eta)$.
The different curvature powers quantify the information supplied by
an exact positive-offset physical family.
\subsection{Geometric setting and observations}
Fix a full-rank lattice $\Lambda\subset\R^2$ and a positive finite number of bounded
$C^\infty$ strictly convex obstacles $D_1,\ldots,D_s$ with everywhere
positive curvature. All distinct translates $D_a+\lambda$, $\lambda\in
\Lambda$, have disjoint closures. Let $A>0$ be the free area per lattice
cell. The unit-speed billiard is prepared with normalized phase volume
$\dd q\dd\vartheta/(2\pi A)$ on the torus complement. We observe the
collision count $N_t$ and, when specified, the first and last collision
positions. Finite horizon, congruence of the obstacles, and symmetry
are not assumed for the general forward and boundary-law results.

Consider a fixed such configuration and a smooth finite-dimensional
family $D_{a,\xi}$ through it. The parameter $\xi$ ranges over a compact
set $K_0$ in a sufficiently small neighborhood of the reference parameter.
The general geometric conclusions are local about an arbitrary reference configuration,
not only about circles. A compact family is covered by finitely many
such neighborhoods, with common constants after taking finite extrema.
Boundaries may be represented by smooth support functions. At every
fixed derivative order, constants can depend on the corresponding
higher norms of this family, on its positive separation and curvature
bounds, and on the lattice. They do not depend on the collision order.

Let $\mathcal E$ be the finite set of oriented closest pairs attaining
the minimum gap in the reference configuration, modulo common lattice
translation. An index $e=(a,b,\lambda)$ represents a segment from
$D_{a,\xi}$ to $D_{b,\xi}+\lambda$; reversing it gives
$\bar e=(b,a,-\lambda)$. Section~\ref{sec:g-localization} proves that
these pairs continue smoothly, remain unobstructed, and include every
channel contributing to a common ground-onset collar. Write
\begin{equation}\label{eq:g-general-data}
 \begin{split}
 g_*&=\min_{e\in\mathcal E}g_e,\qquad
 c_{e,b}=1+g_e\kappa_{e,b}\quad(b=0,1),\\
 c_e&=\sqrt{c_{e,0}c_{e,1}},\qquad \gamma_e=\arcosh c_e,
 \qquad d_{j,e}=t-jg_e.
 \end{split}
\end{equation}
Here $\kappa_{e,0}$ and $\kappa_{e,1}$ are the positive curvatures at
its two contact points. They need not be equal. A channel event
$E_{j,e}(d)$ consists of the $j+1$ impacts alternating between the
specified facing patches in the window $t=jg_e+d$. These patches are
fixed around the reference contacts and chosen before observation.
They identify a measurable part of a physical collision record, not an
extra latent observation. At a nonminimal channel's own onset this
selection is required; the unselected count law there is a different
observation.

Parameter derivatives are multi-index derivatives in a fixed Euclidean
parameter chart; offset derivatives hold $d$ fixed as an independent
coordinate. Matrix norms are Euclidean operator norms unless a trace
norm is explicitly indicated. The growing record arrays later use the
maximum norm. These conventions distinguish uniform normalized
amplitudes from physical-time derivatives, which can carry powers of $j$.


\subsection{The nonlinear relative law}
Fix an oriented channel $e$ and abbreviate $g=g_e$, $\gamma=\gamma_e$.
Use transverse endpoint coordinates $u,v$ at its two contacts. Let
$W_j(u,v)$ be the stationary length of its $j$ alternating flights,
$E_j=W_j-jg$, $d_j^0=-W_{j,uv}(0,0)$, and $b_j=-W_{j,uv}/d_j^0$.
For $p=j\bmod2$ put
\[
 a_b=\frac{\sqrt{c_0c_1}\sinh\gamma}{g c_{1-b}},\qquad
 q_p=\sqrt{a_0a_p},\qquad d_j^0=\frac{q_p}{\sinh(j\gamma)}.
\]
The stationary bridge and these formulas are proved in
Section~\ref{sec:g-action}. All norms in the next theorem are on fixed
common neighborhoods and include the indicated geometric derivatives.

\begin{theorem}[Relative boundary law]\label{thm:v8-main-relative}
There are positive smooth functions $B_b$ and strictly convex smooth
functions $S_b$, $b=0,1$, with
\[
 S_b(0)=S_b'(0)=0,\qquad S_b''(0)=a_b,\qquad B_b(0)=1,
\]
a common endpoint box, a number $d_*>0$, and $0<\tau<1$ such that,
for every fixed derivative order $k$ and every $j\ge1$,
\begin{equation}\label{eq:v8-relative}
 \|E_j-S_0\oplus S_p\|_{C^k}
 +\|b_j-B_0\otimes B_p\|_{C^k}\le C_k\tau^j,
                    \qquad p=j\bmod2.
\end{equation}
The action difference vanishes to second order at the origin. Define
\begin{equation}\label{eq:v8-boundary-integral}
 \mathcal F_p(d)=\frac{q_p}{\pi d^2}
  \int(d-S_0(u)-S_p(v))_+B_0(u)B_p(v)\,\dd u\dd v.
\end{equation}
Then $\mathcal F_p$ extends smoothly to $[0,d_*]$, is positive there,
and has value one at zero. The exact selected-channel probabilities obey
\begin{equation}\label{eq:v8-relative-integral}
 \left\|\frac{2A\sinh(j\gamma)}{d^2}
           \Pr(E_{j,e}(d))-\mathcal F_p(d)\right\|_{C^k([0,d_*])}
                                  \le C_k\tau^j.
\end{equation}
The norm includes every fixed mixed geometric and offset derivative.
The limiting conditional endpoint--residual law at $d>0$ has density
proportional to
\[
 B_0(u)B_p(v)\one_{\{r>0,\ S_0(u)+S_p(v)+r<d\}}.
\]
All assertions hold uniformly on the compact local families specified
above. For analytic families the endpoint constructions are analytic.
\end{theorem}

\begin{proof}[Proof and dependencies]
Lemma~\ref{lem:g-relative} constructs the finite bridge and its relative
determinant. Lemma~\ref{lem:v4-halfline} constructs $S_b,B_b$ with their
normalizations. Theorem~\ref{thm:v4-factorization} proves
\eqref{eq:v8-relative} by comparing two separated perturbation blocks in
trace norm, including their mixed derivatives. Theorem~\ref{thm:v4-law}
proves \eqref{eq:v8-relative-integral} using the common Morse domain of
Lemma~\ref{lem:g-radial}, and identifies the conditional density.
These proofs are supplied in full in the following sections.
\end{proof}

At fixed positive $d$, the boundary functions in this theorem retain
nonlinear contact information. Proposition~\ref{prop:v4-quartic} computes
a nonzero variation of the first nonlinear coefficient, and
Theorem~\ref{thm:v4-jet-fiber} realizes it with fixed area and unchanged
leading data for the selected horizontal hierarchy. Higher-order
realizations and complete contact-jet proofs are given in
Theorem~\ref{thm:v5-realization} and the subsequent appendix. These are
specific demonstrations that the boundary law is not determined by the
linear multiplier. They concern the stated selected hierarchy, not all
channels simultaneously.

The argument uses familiar ingredients for a different target estimate.
Hill identities relate action Hessians and monodromy \cite{Hill}; here
the determinant is compared relatively after nonlinear boundary
perturbations, then integrated against physical residual time. Spectral
and marked-length inverse results \cite{Zelditch,DSKL,FinamoreLeguil}
use different data and hypotheses. The proposition-level operator comparison in
Section~\ref{sec:v9-operators} separates the classical identities,
the geometric trace-norm estimate, and stable physical integration.
Section~\ref{sec:v10-deautoconvolution} gives the closest inverse
autoconvolution comparison, including the singular kernel and observation
topology. Section~\ref{sec:v5-comparison} compares the inverse observations;
neither discussion makes an exhaustive priority claim.

\subsection{Organization and observation boundaries}
The geometric and relative-flux proofs precede their experimental
consequences. The nonlinear compatibility and profile inverse follow the
operator comparison. Full-profile acquisition and endpoint testing then
form the observation part of the article. The
finite-flight exact-family inverse and its acquisition theorems form
a separate sequence. Smooth remainder robustness is stated as
a separate experiment. Complete contact-jet, conditional-position,
coalescence, calibration, record-response, and circular proofs are
collected in the appendices; no additional hypothesis on the general
forward theorem is introduced by those applications.

Long bridges supply uniform physical experiments and a compatible
pair of nonlinear boundary profiles. They do not supply otherwise
inaccessible full contact coordinates: the identical-even
contact germ is already identifiable from a one-flight germ, and the
count inverse below uses only flight numbers $1,2,3$. The exact-family
sampling theorem counts independent full-phase preparations and retains
only binary reports. The endpoint experiment retains more data. None of
the count lower bounds applies to unrestricted collision values,
positions, unbounded flight numbers, or windows outside its declared
collar. Analytic germ identification is also distinct from noise-stable
global continuation.

\subsection{A dependency map}
The main logical branches are as follows. The hypotheses in one
application are not silently transferred to the general forward law.
\begin{center}
\begin{tabular}{@{}p{0.24\textwidth}p{0.31\textwidth}p{0.37\textwidth}@{}}
\toprule
Conclusion & Observation and extra information & Target and control\\
\midrule
Relative physical law & Selected physical channel, general unequal contacts &
Uniform relative $C^k$ action, flux and physical law from localized
determinants and residual integration\\[3pt]
Nonlinear profile inverse & Two oriented even laws; known $A,\gamma,g$, patches;
$C^3$ data norm & Two full energy profiles: $C^3$ law error to $C^0$ profile error;
relative law and Volterra inversion\\[3pt]
Full-profile acquisition & Selected-event binary reports; same known
quantities; supplied regularity and relative bounds & Uniform error of both full profiles; positive-node regularization
and fully charged preparations\\[3pt]
Endpoint testing & Endpoints and residual time; joint tangent regime &
Total variation and bounded risks; relative transfer and quadratic
overlap\\[3pt]
Contact-graph jets & Selected coefficient germ; identical even facing
graphs & Fixed-order graph jets in coefficient norm; triangular inverse;
one-flight input sufficient\\[3pt]
Exact count inverse & Four unlabelled fixed-window means; exact analytic family; unknown gap &
Gap loss $\delta$ and curvature-multiset loss $\varepsilon$;
sharp joint preparation order in its bounded-flight model\\[3pt]
Smooth-envelope optimum & Three binary indicators; leading family;
unknown $C^m$ remainders and strict slack; not an exact-law oracle & Same two finite-dimensional losses; different sharp preparation
order from nuisance splicing and entropy\\
\bottomrule
\end{tabular}
\end{center}

In the first two profile problems the geometric normalization and
orientation labels are supplied, and the target is a pair of functions.
In the exact-family problem four unlabelled means recover finite
symmetric shape coordinates and the gap in that family; its self-calibration
does not remove supplied information from the profile experiment.
The smooth-envelope problem retains a leading family but admits free
nuisance functions. Its minimax alternatives need not be billiard tables.
The respective uniform-profile, curvature-multiset and gap losses are
not interchangeable. These information sets are part of the statements,
not conventions inferred from the choice of reconstruction method.

The full-profile assertions are Theorems~\ref{thm:v9-compatibility},
\ref{thm:v9-profile-stability} and~\ref{thm:v10-acquisition}.
The exact-family benchmark is Theorems~\ref{thm:v8-four}
and~\ref{thm:v8-joint-minimax}; its hypotheses and losses differ from
those of the smooth-envelope benchmark,
Theorem~\ref{thm:v9-envelope-minimax}. None of these comparisons
transfers the information supplied in one experiment to another.
